% Options for packages loaded elsewhere
\PassOptionsToPackage{unicode}{hyperref}
\PassOptionsToPackage{hyphens}{url}
\documentclass[
  11pt,
]{article}
\usepackage{xcolor}
\usepackage[margin=2.2cm]{geometry}
\usepackage{amsmath,amssymb}
\setcounter{secnumdepth}{-\maxdimen} % remove section numbering
\usepackage{iftex}
\ifPDFTeX
  \usepackage[T1]{fontenc}
  \usepackage[utf8]{inputenc}
  \usepackage{textcomp} % provide euro and other symbols
\else % if luatex or xetex
  \usepackage{unicode-math} % this also loads fontspec
  \defaultfontfeatures{Scale=MatchLowercase}
  \defaultfontfeatures[\rmfamily]{Ligatures=TeX,Scale=1}
\fi
\usepackage{lmodern}
\ifPDFTeX\else
  % xetex/luatex font selection
  \setmainfont[]{DejaVu Serif}
  \setsansfont[]{DejaVu Sans}
  \setmonofont[]{DejaVu Sans Mono}
\fi
% Use upquote if available, for straight quotes in verbatim environments
\IfFileExists{upquote.sty}{\usepackage{upquote}}{}
\IfFileExists{microtype.sty}{% use microtype if available
  \usepackage[]{microtype}
  \UseMicrotypeSet[protrusion]{basicmath} % disable protrusion for tt fonts
}{}
\makeatletter
\@ifundefined{KOMAClassName}{% if non-KOMA class
  \IfFileExists{parskip.sty}{%
    \usepackage{parskip}
  }{% else
    \setlength{\parindent}{0pt}
    \setlength{\parskip}{6pt plus 2pt minus 1pt}}
}{% if KOMA class
  \KOMAoptions{parskip=half}}
\makeatother
\setlength{\emergencystretch}{3em} % prevent overfull lines
\providecommand{\tightlist}{%
  \setlength{\itemsep}{0pt}\setlength{\parskip}{0pt}}
% ---- custom preamble for the RoadFury -> SE2RoadNet speaker script ----
\usepackage{xcolor}
\usepackage{framed}
\usepackage{titlesec}
\usepackage{enumitem}
\usepackage{parskip}

\definecolor{deckteal}{HTML}{16242A}
\definecolor{deckorange}{HTML}{E0772B}
\definecolor{deckslate}{HTML}{45595B}

\setcounter{secnumdepth}{0}
\linespread{1.05}

% coloured, unnumbered headings matching the deck palette
\titleformat{\section}{\normalfont\Large\bfseries\color{deckorange}}{}{0pt}{}
\titleformat{\subsection}{\normalfont\large\bfseries\color{deckteal}}{}{0pt}{}
\titlespacing*{\section}{0pt}{1.5em}{0.55em}
\titlespacing*{\subsection}{0pt}{1.05em}{0.35em}

% markdown blockquotes ( "> VERIFY: ..." ) -> small italic left-bar callout
\renewenvironment{quote}
  {\par\smallskip\noindent
   \def\FrameCommand{{\color{deckorange}\vrule width 2.5pt}\hspace{10pt}}%
   \MakeFramed{\advance\hsize-\width\FrameRestore}%
   \footnotesize\itshape\color{deckslate}}
  {\par\endMakeFramed\smallskip}

\setlist{topsep=2pt,parsep=1pt,itemsep=1pt,leftmargin=1.4em}

% DejaVu Serif lacks U+2713/U+2717; borrow the check/cross from DejaVu Sans
\usepackage{newunicodechar}
\newfontfamily\dvsans{DejaVu Sans}
\newunicodechar{✓}{{\dvsans ✓}}
\newunicodechar{✗}{{\dvsans ✗}}
\usepackage{bookmark}
\IfFileExists{xurl.sty}{\usepackage{xurl}}{} % add URL line breaks if available
\urlstyle{same}
% fallback for those not using the hyperref driver hyperxmp:
\makeatletter
\@ifundefined{xmpquote}{\newcommand{\xmpquote}[1]{#1}}{}
\makeatother
\hypersetup{
  pdftitle={RoadFury → SE2RoadNet --- Lời thuyết trình cho thầy},
  pdfauthor={Trần Chí Nguyên · Huỳnh Trung Kiệt · Đào Sỹ Duy Minh},
  hidelinks,
  pdfcreator={LaTeX via pandoc}}

\title{RoadFury → SE2RoadNet --- Lời thuyết trình cho thầy}
\usepackage{etoolbox}
\makeatletter
\providecommand{\subtitle}[1]{% add subtitle to \maketitle
  \apptocmd{\@title}{\par {\large #1 \par}}{}{}
}
\makeatother
\subtitle{Bản giải thích phương pháp · The Self-Driving Car Testing
Competition · ICST 2026}
\author{Trần Chí Nguyên · Huỳnh Trung Kiệt · Đào Sỹ Duy Minh}
\date{Văn nói (học trò → thầy) · deck 37 slide (s00--s36) · biên dịch
XeLaTeX}

\begin{document}
\maketitle

{
\setcounter{tocdepth}{1}
\tableofcontents
}
\section{Mở đầu}\label{mux1edf-ux111ux1ea7u}

\subsection{Cách nói mở đầu (trước khi bấm slide) ---
{[}20s{]}}\label{cuxe1ch-nuxf3i-mux1edf-ux111ux1ea7u-trux1b0ux1edbc-khi-bux1ea5m-slide-20s}

Em chào thầy và cả lớp ạ. Nhóm em xin trình bày đề tài \textbf{ưu tiên
kiểm thử xe tự lái}, tên phương pháp là \textbf{RoadFury tiến tới
SE2RoadNet}. Buổi hôm nay em xin đi khoảng 15 phút, trọng tâm nằm ở phần
\textbf{phương pháp} và phần \textbf{kết quả}.

Nếu thầy cho phép, em xin nói trước \textbf{một câu để thầy nắm ý
chính}: nhóm em xây một mô hình chấm điểm con đường để xếp thứ tự chạy
kiểm thử; điểm mới của nhóm em là mô hình đó \textbf{bất biến với phép
xoay con đường --- bằng chính thiết kế kiến trúc, chứng minh được trên
giấy, chứ không phải bằng thủ thuật tăng cường dữ liệu}. Trong suốt bài,
em sẽ quay lại làm rõ ý này cho thầy.

\subsection{s00 · Trang tiêu đề ---
{[}15s{]}}\label{s00-trang-tiuxeau-ux111ux1ec1-15s}

Đây là đề tài của nhóm em, gồm ba thành viên. Tên phương pháp có hai vế:
\textbf{RoadFury} là mô hình nền tảng nhóm em tự xây trước đó, và
\textbf{SE2RoadNet} là đề xuất chính mà em sẽ trình bày kỹ hôm nay. Mũi
tên giữa hai tên hàm ý: nhóm em không làm lại từ đầu, mà \textbf{vá đúng
hai điểm yếu} của RoadFury để ra SE2RoadNet.

\subsection{s01 · Mục lục ---
{[}20s{]}}\label{s01-mux1ee5c-lux1ee5c-20s}

Bài của em có năm phần. Phần một là bối cảnh và bài toán, để thầy thấy
vì sao việc này khó và đáng làm. Phần hai và ba là \textbf{phương pháp}
--- đây là phần em xin dành nhiều thời gian nhất để thầy hiểu rõ nhóm em
đã làm gì. Phần bốn là \textbf{đánh giá thực nghiệm}. Phần năm em xin
nói thẳng những hạn chế còn lại. Cuối bài em có thêm năm slide dự phòng,
phòng khi thầy hỏi sâu.

\section{Phần 1 --- Bài toán}\label{phux1ea7n-1-buxe0i-touxe1n}

\subsection{s02 · Chuyển phần 1 ---
{[}5s{]}}\label{s02-chuyux1ec3n-phux1ea7n-1-5s}

Thưa thầy, em xin bắt đầu bằng bài toán ạ.

\subsection{s03 · Bối cảnh: kiểm thử xe tự lái trong mô phỏng ---
{[}45s{]}}\label{s03-bux1ed1i-cux1ea3nh-kiux1ec3m-thux1eed-xe-tux1ef1-luxe1i-trong-muxf4-phux1ecfng-45s}

Để kiểm thử phần mềm lái tự động, người ta không cho xe chạy ngoài đường
thật mà chạy trong \textbf{mô phỏng vật lý}, ví dụ nền BeamNG.tech. Mỗi
lần build phần mềm, hệ thống sinh \textbf{hàng nghìn} con đường khác
nhau một cách tự động, rồi cho xe ảo chạy giữ làn trên từng con đường:
xe giữ được làn thì \textbf{PASS}, lệch ra ngoài thì \textbf{FAIL}.

Vấn đề, thưa thầy, là \textbf{mỗi con đường chạy mô phỏng vật lý mất từ
vài giây tới cả phút}; gom hàng nghìn kịch bản lại thì tốn rất nhiều giờ
máy. Mà phần lớn kịch bản lại PASS, chỉ chừng ba tới bốn phần mười là
FAIL. Cho nên thay vì chạy tràn lan theo thứ tự ngẫu nhiên, ý tưởng là
\textbf{xếp những kịch bản dễ FAIL lên chạy trước}, để lộ lỗi sớm và
tiết kiệm chi phí.

\subsection{s04 · Phát biểu bài toán ---
{[}45s{]}}\label{s04-phuxe1t-biux1ec3u-buxe0i-touxe1n-45s}

Em xin hình thức hóa để thầy tiện theo dõi. \textbf{Đầu vào} của mô hình
là một con đường, biểu diễn bằng một chuỗi các điểm hai chiều --- tùy
con đường mà có từ 64 tới 197 điểm. Xin thầy lưu ý: đầu vào \textbf{chỉ
có hình học con đường}, không có bất kỳ dữ liệu hành vi nào của xe.

Mô hình của nhóm em, em gọi là bộ chấm điểm f, nhận con đường đó và trả
về \textbf{một xác suất FAIL trong khoảng 0 tới 1}. \textbf{Đầu ra} cuối
cùng là một thứ tự: xếp các kịch bản theo điểm giảm dần. Mục tiêu là làm
sao \textbf{lỗi bị đẩy lên đầu hàng đợi càng nhiều càng tốt}, và điều đó
được đo bằng chỉ số tên là \textbf{APFD}, em sẽ giải thích ngay sau đây.

\subsection{s05 · Vì sao cần ưu tiên ---
{[}40s{]}}\label{s05-vuxec-sao-cux1ea7n-ux1b0u-tiuxean-40s}

Để thầy hình dung APFD một cách trực quan: hãy tưởng tượng em vẽ đường
cong ``đã chạy bao nhiêu phần trăm test thì phát hiện được bao nhiêu
phần trăm lỗi''. Nếu xếp hạng tốt, lỗi dồn lên đầu, đường cong này
\textbf{dốc đứng lên ngay từ sớm}; nếu xếp hạng kém, lỗi nằm cuối, đường
cong đi ngang mãi rồi mới lên. \textbf{APFD chính là diện tích dưới
đường cong đó} --- càng cao nghĩa là lộ lỗi càng sớm.

Điểm em muốn thầy nhớ ở đây là: bài toán này \textbf{không phải ``chạy
cho nhanh hơn''}, mà là \textbf{``chạy đúng thứ tự''}. Cùng một tập lỗi,
chỉ cần đổi thứ tự chạy là chi phí đã khác nhau rất nhiều.

\subsection{s06 · Bộ dữ liệu SensoDat ---
{[}45s{]}}\label{s06-bux1ed9-dux1eef-liux1ec7u-sensodat-45s}

Nhóm em dùng bộ dữ liệu công khai \textbf{SensoDat}, công bố tại hội
nghị MSR 2024. Đây là các kịch bản xe tự lái sinh tự động trên BeamNG.
Bộ này chia làm ba phần: tập \textbf{huấn luyện} gần 28.800 kịch bản,
tập \textbf{kiểm tra} hơn 7.200 kịch bản cùng phân phối, và một tập
riêng gọi là \textbf{Competition} gồm 956 kịch bản.

Điều quan trọng, thưa thầy, là tập Competition \textbf{lệch phân phối}
so với tập huấn luyện --- thuật ngữ là out-of-distribution. Cụ thể,
đường trong tập này ngắn hơn hẳn, chỉ từ 129 tới 229 mét, trong khi huấn
luyện có đường dài tới hơn 450 mét. Nhóm em cố tình đánh giá trên tập
lệch này để chứng minh mô hình \textbf{tổng quát hóa được}, chứ không
chỉ học thuộc phân phối huấn luyện.

\subsection{s07 · Ví dụ kịch bản: hình dạng quyết định nhãn ---
{[}35s{]}}\label{s07-vuxed-dux1ee5-kux1ecbch-bux1ea3n-huxecnh-dux1ea1ng-quyux1ebft-ux111ux1ecbnh-nhuxe3n-35s}

Đây là dữ liệu thật, thưa thầy. Quan sát của nhóm em: đường
\textbf{PASS} thường cong mượt, còn đường \textbf{FAIL} thường có
\textbf{chicane} --- tức những khúc cua zigzag gắt, độ cong đổi dấu liên
tục làm xe không kịp bám làn.

Em xin nhấn mạnh một điểm, vì nó là nền cho toàn bộ phương pháp:
\textbf{không có tọa độ tuyệt đối nào của con đường quyết định nhãn PASS
hay FAIL --- chỉ có hình dạng quyết định}. Một con đường zigzag thì dù
em đặt nó ở đâu trên bản đồ, xoay theo hướng nào, nó vẫn là con đường
zigzag khó lái. Đây chính là gợi ý cho ràng buộc bất biến ở phần sau.

\subsection{s08 · Công trình liên quan: ba hướng ---
{[}40s{]}}\label{s08-cuxf4ng-truxecnh-liuxean-quan-ba-hux1b0ux1edbng-40s}

Trước nhóm em, có ba hướng tiếp cận. Hướng thứ nhất là \textbf{tìm kiếm
theo độ đa dạng} --- dùng giải thuật di truyền chọn ra tập test đa dạng
hành vi; nhược điểm là \textbf{đa dạng chưa chắc lộ lỗi}. Hướng thứ hai
là \textbf{học máy trên đặc trưng thủ công} --- ví dụ mô hình ITS4SDC
dùng mạng LSTM trên vài đặc trưng chuỗi; rẻ và dễ giải thích, nhưng ít
đặc trưng và \textbf{phụ thuộc khung quy chiếu}. Hướng thứ ba là
\textbf{học sâu} --- mạng đồ thị, ResNet, hay Transformer như RoadFury
của nhóm em; cho APFD cao nhất nhưng vẫn \textbf{không có gì bảo đảm
tính bất biến}.

Khoảng trống mà nhóm em thấy, và cũng là động lực của đề tài:
\textbf{chưa phương pháp nào bảo đảm được tính bất biến với phép xoay
hay với tần số lấy mẫu --- tất cả đều chỉ dựa vào tăng cường dữ liệu,
tức là gần đúng.}

\subsection{s09 · Công trình liên quan: tổng hợp ---
{[}35s{]}}\label{s09-cuxf4ng-truxecnh-liuxean-quan-tux1ed5ng-hux1ee3p-35s}

Bảng này tổng hợp lại. Xin thầy đọc theo hai chiều. Chiều APFD thì các
phương pháp tăng dần từ ngẫu nhiên khoảng 0,49 lên tới các mô hình học
sâu khoảng 0,80. Nhưng chiều thứ hai, cột \textbf{``bất biến''}, thì
\textbf{mọi phương pháp trước đều không có}.

Thông điệp của nhóm em nằm ở đây, thưa thầy: nhóm em \textbf{không cố
chạy đua thêm vài phần nghìn APFD}, mà đóng góp một \textbf{trục bảo đảm
hoàn toàn mới} mà chưa ai có. Đó là tinh thần của cả đề tài.

\section{Phần 2 --- RoadFury (mô hình nền
tảng)}\label{phux1ea7n-2-roadfury-muxf4-huxecnh-nux1ec1n-tux1ea3ng}

\subsection{s10 · Chuyển phần 2 ---
{[}5s{]}}\label{s10-chuyux1ec3n-phux1ea7n-2-5s}

Phần hai, em xin giới thiệu RoadFury --- mô hình nền tảng của nhóm em ạ.

\subsection{s11 · RoadFury: sơ đồ tổng thể ---
{[}30s{]}}\label{s11-roadfury-sux1a1-ux111ux1ed3-tux1ed5ng-thux1ec3-30s}

RoadFury hoạt động theo bốn bước. Thứ nhất, chuẩn hóa con đường về 197
điểm và \textbf{trích ra 10 kênh đặc trưng} tại mỗi điểm. Thứ hai, đưa
vào một \textbf{Transformer} để đọc con đường như một chuỗi. Thứ ba,
dùng kỹ thuật \textbf{SWA} --- trung bình trọng số các mô hình cuối quá
trình huấn luyện --- để mô hình ổn định hơn. Thứ tư, lúc suy luận thì
chấm điểm và sắp xếp giảm dần. Đây là một mô hình rất mạnh, nhưng nó có
\textbf{hai điểm mù} mà em sẽ chỉ ra ngay ở slide sau --- và chính hai
điểm mù này sinh ra SE2RoadNet.

\subsection{s12 · RoadFury chi tiết và hai điểm mù ---
{[}55s{]}}\label{s12-roadfury-chi-tiux1ebft-vuxe0-hai-ux111iux1ec3m-muxf9-55s}

Em xin đi vào chi tiết để thầy thấy rõ hai điểm mù. Bên trái là 10 kênh
đặc trưng: chiều dài đoạn, biến thiên góc, độ cong, độ giật của độ cong,
và nhiều đại lượng khác. \textbf{Nhưng có ba kênh --- sin của góc, cos
của góc, và góc tuyệt đối --- là phụ thuộc khung quy chiếu.} Nghĩa là
gì, thưa thầy? Nghĩa là nếu em xoay con đường đi một góc, ba kênh này
đổi giá trị ngay, dù con đường thật ra vẫn y hệt.

Bên phải là kiến trúc Transformer, khoảng 829 nghìn tham số, và nó dùng
\textbf{mã hóa vị trí tuyệt đối} --- đây là điểm mù thứ hai. Kết quả thì
rất tốt: APFD khoảng 0,80, AUC 0,917, đứng đầu cuộc thi. \textbf{Nhưng
điểm cao mà không có bảo đảm}: chỉ cần xoay con đường là điểm số trôi
đi. Với một hệ thống an toàn như xe tự lái, đó là một lỗ hổng nhóm em
không thể bỏ qua.

\subsection{s13 · Cầu nối: vá hai điểm mù ---
{[}30s{]}}\label{s13-cux1ea7u-nux1ed1i-vuxe1-hai-ux111iux1ec3m-muxf9-30s}

Vậy nên toàn bộ SE2RoadNet, thưa thầy, chỉ là \textbf{vá đúng hai điểm
mù đó, có chủ đích}. Điểm mù thứ nhất --- nhạy với phép xoay --- nhóm em
vá bằng cách \textbf{bỏ ba kênh phụ thuộc khung, chỉ giữ 7 kênh bất
biến}. Điểm mù thứ hai --- nhạy với cách lấy mẫu --- nhóm em vá bằng
cách \textbf{thay mã hóa vị trí tuyệt đối bằng một cơ chế chỉ dựa trên
khoảng cách dọc đường}.

Tinh thần chung là chuyển từ \textbf{tăng cường dữ liệu, vốn chỉ gần
đúng}, sang \textbf{ràng buộc bằng kiến trúc, vốn có bảo đảm}. Đây là ý
em mong thầy nắm nhất trước khi vào phần ba.

\section{Phần 3 --- SE2RoadNet (đề xuất
chính)}\label{phux1ea7n-3-se2roadnet-ux111ux1ec1-xuux1ea5t-chuxednh}

\subsection{s14 · Chuyển phần 3 ---
{[}5s{]}}\label{s14-chuyux1ec3n-phux1ea7n-3-5s}

Và đây là phần trọng tâm ạ: SE2RoadNet.

\subsection{s15 · SE2RoadNet: tổng quan ---
{[}30s{]}}\label{s15-se2roadnet-tux1ed5ng-quan-30s}

Sơ đồ tổng thể gồm: 7 kênh đặc trưng bất biến, đưa qua 6 khối mà nhóm em
gọi là \textbf{InvariantBlock}, rồi gộp lại và cho ra điểm FAIL. Em xin
nói luôn \textbf{kết quả then chốt} để thầy có đích theo dõi: khi nhóm
em xoay con đường bằng nhiều góc khác nhau rồi chấm lại, độ chênh APFD
là \textbf{đúng bằng 0} --- không phải xấp xỉ 0, mà bằng 0. Bây giờ em
xin giải thích bằng bốn bước làm sao đạt được điều đó.

\subsection{s16 · Ý tưởng cốt lõi: bất biến SE(2) ---
{[}50s{]}}\label{s16-uxfd-tux1b0ux1edfng-cux1ed1t-luxf5i-bux1ea5t-biux1ebfn-se2-50s}

Trước hết, xin thầy cho em giải thích \textbf{SE(2) là gì}. SE(2) là
nhóm các \textbf{phép biến đổi cứng của mặt phẳng} --- gồm phép
\textbf{xoay} và phép \textbf{tịnh tiến}. Nói nôm na, thưa thầy: nếu em
cầm bản vẽ con đường lên rồi xoay đi một góc bất kỳ, hoặc dịch sang chỗ
khác, hoặc đặt lại gốc tọa độ --- thì đó \textbf{vẫn là đúng con đường
đó}, chiếc xe vẫn lái y hệt, PASS vẫn PASS và FAIL vẫn FAIL.

Cho nên điểm số mà mô hình chấm \textbf{bắt buộc phải giống hệt nhau}
trong cả bốn trường hợp. Trong slide có bốn hình cùng một con đường, và
cả bốn đều cho cùng một điểm. Em xin nhấn mạnh sự khác biệt cốt lõi:
\textbf{tăng cường dữ liệu chỉ làm mô hình quen với phép xoay; còn kiến
trúc của nhóm em thì không cho phép điểm số đổi --- về mặt toán học là
bất khả.}

\begin{quote}
Chốt ý cho thầy: bất biến ở đây là \textbf{tính chất của kiến trúc},
đúng theo định nghĩa, không phụ thuộc vào việc mô hình được huấn luyện
tốt hay xấu.
\end{quote}

\subsection{s17 · Bước 1: bảy kênh đặc trưng bất biến ---
{[}55s{]}}\label{s17-bux1b0ux1edbc-1-bux1ea3y-kuxeanh-ux111ux1eb7c-trux1b0ng-bux1ea5t-biux1ebfn-55s}

Bước một là chọn đặc trưng. Nhóm em chỉ giữ lại \textbf{7 đại lượng nội
tại} của con đường: chiều dài mỗi đoạn, biến thiên góc, \textbf{độ
cong}, đạo hàm bậc một và bậc hai của độ cong theo chiều dài, chiều dài
cung chuẩn hóa, và độ lệch chuẩn cục bộ của độ cong. Điểm mấu chốt:
\textbf{không còn kênh nào phụ thuộc khung quy chiếu} --- nhóm em đã bỏ
hẳn sin, cos và góc tuyệt đối.

Thầy có thể hỏi: bỏ tọa độ và hướng đi như vậy thì có mất thông tin
không? Câu trả lời dựa trên định lý \textbf{Frenet--Serret}: một đường
cong phẳng được xác định \textbf{duy nhất, sai khác đúng một phép
xoay-tịnh tiến, bởi hàm độ cong của nó}. Nói cách khác, nếu em biết con
đường ``cong bao nhiêu tại mỗi điểm dọc theo chiều dài'', thì em đã nắm
trọn hình dạng của nó rồi. Cho nên bỏ ba kênh kia là \textbf{không mất
thông tin gì cả} --- chỉ bỏ đúng phần dư thừa mà lại chính là phần gây
nhạy cảm với phép xoay.

\begin{quote}
Chốt ý cho thầy: tính bất biến được bảo đảm \textbf{ngay từ khâu đặc
trưng, trước khi mô hình học một tham số nào}.
\end{quote}

\subsection{s18 · Bước 2: kiến trúc InvariantBlock ---
{[}45s{]}}\label{s18-bux1b0ux1edbc-2-kiux1ebfn-truxfac-invariantblock-45s}

Bước hai là kiến trúc. Mỗi điểm trên đường được biến thành một vector
192 chiều --- gọi là một token. Nhóm em thêm một token tổng hợp đặc
biệt, rồi cho qua \textbf{6 khối InvariantBlock}. Mỗi khối là một lớp
\textbf{chú ý đa đầu} tám đầu, cộng mạng truyền thẳng và chuẩn hóa, kèm
theo một thành phần đặc biệt em sẽ nói ở bước ba. Cuối cùng gộp thông
tin về token tổng hợp và cho ra một con số là xác suất FAIL.

Về quy mô: tổng cộng \textbf{2,11 triệu tham số}, gấp khoảng hai lần
rưỡi RoadFury, nhưng vẫn nhẹ --- huấn luyện chỉ \textbf{24 phút}. Và xin
thầy lưu ý: nhóm em dùng \textbf{một cấu hình duy nhất}, không tinh
chỉnh riêng cho từng bộ dữ liệu.

\subsection{s19 · Bước 3: chú ý theo hiệu chiều dài cung ---
{[}60s{]}}\label{s19-bux1b0ux1edbc-3-chuxfa-uxfd-theo-hiux1ec7u-chiux1ec1u-duxe0i-cung-60s}

Bước ba là phần kỹ thuật lõi, em xin nói kỹ ạ. Transformer cần biết thứ
tự các điểm, nên thông thường người ta cộng thêm \textbf{mã hóa vị trí}
--- một tín hiệu theo \textbf{chỉ số} của điểm: điểm thứ nhất, thứ hai,
thứ ba. Vấn đề, thưa thầy, là chỉ số tuyệt đối này \textbf{phụ thuộc vào
cách mình lấy mẫu con đường}: lấy 64 điểm hay 197 điểm thì cùng một vị
trí vật lý lại rơi vào chỉ số khác nhau, và thế là tính bất biến bị phá
vỡ.

Giải pháp của nhóm em: thay vì mã hóa vị trí tuyệt đối, nhóm em thêm một
\textbf{độ lệch --- gọi là bias --- vào ma trận chú ý}, và bias này
\textbf{chỉ phụ thuộc vào hiệu chiều dài cung giữa hai điểm}, tức khoảng
cách dọc theo con đường giữa chúng. Hiệu số này là đại lượng nội tại ---
xoay hay dịch con đường đều không làm nó đổi --- nên tính bất biến được
giữ trọn vẹn. Về mặt cài đặt, nhóm em đưa hiệu số đó qua một hàm với 32
tần số ngẫu nhiên cố định rồi qua một mạng nhỏ.

Cái giá phải trả là chi phí tính toán bậc hai theo số điểm, nên huấn
luyện mất 24 phút; hướng cải tiến sắp tới của nhóm em là thay bằng một
kỹ thuật rẻ hơn tên là RoPE.

\begin{quote}
Chốt ý cho thầy: mã hóa vị trí tuyệt đối là điểm mù thứ hai của
RoadFury; nhóm em thay nó bằng thứ \textbf{chỉ nhìn vào khoảng cách
tương đối dọc đường}, nhờ vậy giữ được bất biến.
\end{quote}

\subsection{s20 · Bước 4: huấn luyện ---
{[}45s{]}}\label{s20-bux1b0ux1edbc-4-huux1ea5n-luyux1ec7n-45s}

Bước bốn là huấn luyện. Điểm đáng nói là \textbf{hàm mất mát Focal}. Vì
FAIL chỉ chiếm ba tới bốn phần mười, nếu dùng hàm mất mát thường thì mô
hình dễ ``đoán PASS hết'' cho chắc ăn. Hàm Focal thêm một hệ số
\textbf{giảm trọng số các ca dễ và tăng trọng số các ca khó}, buộc mô
hình tập trung vào những ca thật sự khó phân biệt.

Nhóm em có \textbf{quét tham số Focal từ 1 đến 5} và thấy APFD
\textbf{gần như phẳng --- mọi giá trị đều nằm trong khoảng sai số} ---
chứng tỏ mô hình bền vững với tham số này; nhóm em chọn giá trị
\textbf{1,0}. Phần còn lại là các cấu hình huấn luyện tiêu chuẩn: bộ tối
ưu AdamW, lịch học có khởi động, SWA ở giai đoạn cuối. Tất cả gói gọn
trong 24 phút, một cấu hình cho mọi bộ dữ liệu.

\section{Phần 4 --- Đánh giá thực
nghiệm}\label{phux1ea7n-4-ux111uxe1nh-giuxe1-thux1ef1c-nghiux1ec7m}

\subsection{s21 · Chuyển phần 4 ---
{[}5s{]}}\label{s21-chuyux1ec3n-phux1ea7n-4-5s}

Phần bốn, em xin trình bày kết quả ạ.

\subsection{s22 · Chỉ số APFD và giao thức đánh giá ---
{[}45s{]}}\label{s22-chux1ec9-sux1ed1-apfd-vuxe0-giao-thux1ee9c-ux111uxe1nh-giuxe1-45s}

Trước hết là chỉ số. \textbf{APFD} nằm trong khoảng 0 tới 1, càng cao
càng tốt; ngẫu nhiên rơi vào khoảng 0,5. Nhóm em đánh giá theo
\textbf{ba lớp} để chắc chắn, thưa thầy. Lớp một, \textbf{chấm một lượt}
trên toàn bộ 956 kịch bản, được APFD 0,8047. Lớp hai, \textbf{chấm 30
lần}, mỗi lần lấy ngẫu nhiên một phần ba tập, để loại yếu tố may rủi do
thứ tự cố định --- kết quả 0,8048 với độ lệch chuẩn 0,012. Lớp ba là
phép thử em tâm đắc nhất: \textbf{thử phép xoay} --- xoay cả tập bằng
nhiều góc rồi chấm lại, đo độ chênh APFD.

\subsection{s23 · Kết quả headline: bất biến phép xoay bằng 0 ---
{[}70s{]}}\label{s23-kux1ebft-quux1ea3-headline-bux1ea5t-biux1ebfn-phuxe9p-xoay-bux1eb1ng-0-70s}

Đây là kết quả em muốn thầy chú ý nhất trong cả bài. Nhóm em xoay tập
kiểm thử bằng sáu góc khác nhau --- 0, 30, 60, 90, 180 và âm 45 độ ---
rồi chấm lại. Với SE2RoadNet, cả sáu góc cho \textbf{cùng một con số
0,8047}, độ chênh \textbf{đúng bằng 0,0000}. Để so sánh, một mô hình
baseline khi bị xoay thì tụt điểm rõ rệt, độ chênh tới 0,057.

Em xin làm rõ để thầy thấy đây không phải nói quá: \textbf{con số này
không phải ``nằm trong sai số nhỏ'', mà bằng nhau đến từng bit dấu phẩy
động}. Lý do là vì bảy kênh đặc trưng của nhóm em trả về \textbf{đúng
cùng một vector} sau khi xoay, và mạng đã ở chế độ suy luận nên hoàn
toàn tất định, cho nên đầu ra giống hệt. Đây là một trong số rất ít
trường hợp trong học máy mà nhóm em có thể nói \textbf{lý thuyết được
kiểm chứng bằng thực nghiệm, không cần kèm chữ ``xấp xỉ''}.

\subsection{s24 · Bảng xếp hạng ---
{[}55s{]}}\label{s24-bux1ea3ng-xux1ebfp-hux1ea1ng-55s}

Đây là bảng xếp hạng so với các baseline. Em xin trình bày \textbf{trung
thực}, thưa thầy: về APFD thuần, SE2RoadNet ở mức 0,805, tức
\textbf{ngang bằng baseline tốt nhất trong phạm vi sai số} --- nhóm em
\textbf{không khẳng định thắng về APFD}. Nhưng điểm được là: \textbf{AUC
tăng rõ rệt từ 0,917 lên 0,934}, \textbf{và} SE2RoadNet là
\textbf{phương pháp duy nhất có bảo đảm bất biến bằng 0}.

Nói theo ngôn ngữ tối ưu, đây là một \textbf{cải thiện Pareto}: không
thua ở chiều nào, tốt hơn ở AUC, và có thêm một bảo đảm lý thuyết mà
không mô hình nào khác có. Trong bối cảnh an toàn của xe tự lái, em cho
rằng \textbf{một bảo đảm chắc chắn đáng giá hơn vài phần nghìn APFD}.

\subsection{s25 · AUC và APFD: vì sao tách bạch ---
{[}40s{]}}\label{s25-auc-vuxe0-apfd-vuxec-sao-tuxe1ch-bux1ea1ch-40s}

Một điểm em muốn giải thích để thầy khỏi thắc mắc: tại sao nhóm em tách
riêng AUC và APFD. \textbf{AUC} đo khả năng phân loại đúng PASS với FAIL
trên từng cặp. \textbf{APFD} đo tốc độ lộ lỗi theo thứ tự chạy --- đây
mới là thứ bài toán thực sự cần. Trong gần như mọi thí nghiệm của nhóm
em, hai chỉ số này \textbf{không đi cùng nhau}: ở đây AUC tăng nhưng
APFD gần như phẳng. Cho nên khi báo cáo, nhóm em luôn nói rõ đang tối ưu
chỉ số nào, chứ không gộp thành một câu ``tốt hơn'' chung chung.

\section{Phần 5 --- Hạn chế và hướng phát
triển}\label{phux1ea7n-5-hux1ea1n-chux1ebf-vuxe0-hux1b0ux1edbng-phuxe1t-triux1ec3n}

\subsection{s26 · Chuyển phần 5 ---
{[}5s{]}}\label{s26-chuyux1ec3n-phux1ea7n-5-5s}

Cuối cùng, em xin nói thẳng những hạn chế còn lại ạ.

\subsection{s27 · Hạn chế ---
{[}50s{]}}\label{s27-hux1ea1n-chux1ebf-50s}

Nhóm em xin trung thực với thầy. Thứ nhất, như vừa nói, AUC và APFD phân
kỳ, nên phải lập luận cẩn thận metric nào quan trọng khi nào. Thứ hai,
một số cải tiến phụ chỉ giúp \textbf{giảm độ dao động} chứ chưa nâng
được APFD trung bình --- nhóm em coi đó là đóng góp về ``độ ổn định'',
không thổi phồng thành kết quả chính. Thứ ba, và quan trọng nhất về mặt
lý thuyết: \textbf{bất biến phép xoay là chính xác tuyệt đối, nhưng bất
biến với tần số lấy mẫu thì mới chỉ gần đúng} --- độ chênh khoảng 0,001,
rất nhỏ nhưng chưa bằng 0. Ngoài ra, phần chặn an toàn bằng conformal
thì nhóm em vẫn đang hoàn thiện. Em xin báo cáo thẳng thay vì che đi.

\subsection{s28 · Hướng phát triển ---
{[}40s{]}}\label{s28-hux1b0ux1edbng-phuxe1t-triux1ec3n-40s}

Từ những hạn chế đó, nhóm em có bốn hướng đi tiếp. Một, \textbf{tổng
quát hóa lên hơn tám bộ benchmark công khai} bằng đúng một công thức,
không tinh chỉnh riêng. Hai, hoàn thiện phần \textbf{chặn an toàn
conformal} cho vừa đúng vừa hữu ích. Ba, làm cho \textbf{bất biến lấy
mẫu cũng chính xác tuyệt đối} bằng kỹ thuật RoPE thay cho cách hiện tại.
Bốn, dùng \textbf{học tự giám sát gắn với vật lý} để học biểu diễn nền
tốt hơn. Mục tiêu xa là biến một mô hình bất biến cho xe tự lái thành
\textbf{một công thức chung, bất biến và kiểm toán được, cho mọi
benchmark}.

\subsection{s29 · Tài liệu tham khảo ---
{[}15s{]}}\label{s29-tuxe0i-liux1ec7u-tham-khux1ea3o-15s}

Đây là các tài liệu tham khảo. Nguồn dữ liệu là SensoDat; nền tảng kỹ
thuật gồm cơ chế chú ý, hàm Focal, kỹ thuật SWA, và lý thuyết bất biến
nhóm. RoadFury và SE2RoadNet là phương pháp của nhóm em; các baseline
khác nhóm em tái lập trong cùng một khung đánh giá.

\subsection{s30 · Video demo --- {[}20s{]}}\label{s30-video-demo-20s}

Đây là video demo trình bày phương pháp của nhóm em, có link Google
Drive và Facebook ạ. Nếu thầy muốn, em có thể mở phát trực tiếp.

\subsection{s31 · Cảm ơn --- {[}10s{]}}\label{s31-cux1ea3m-ux1a1n-10s}

Phần trình bày của nhóm em đến đây là hết. Em xin cảm ơn thầy và cả lớp
đã lắng nghe, và em xin sẵn sàng nhận câu hỏi ạ.

\section{Dự phòng --- trả lời khi thầy
hỏi}\label{dux1ef1-phuxf2ng-trux1ea3-lux1eddi-khi-thux1ea7y-hux1ecfi}

\subsection{s32 · Dự phòng B1: APFD tính tay --- {[}khi được
hỏi{]}}\label{s32-dux1ef1-phuxf2ng-b1-apfd-tuxednh-tay-khi-ux111ux1b0ux1ee3c-hux1ecfi}

Nếu thầy muốn thấy APFD tính cụ thể: giả sử 5 test, 2 cái FAIL. Nếu xếp
hai cái FAIL lên vị trí 1 và 2 thì APFD bằng 0,80; nếu xếp xuống cuối,
vị trí 4 và 5, thì APFD chỉ còn 0,20. Cùng một tập lỗi, chỉ khác thứ tự,
mà điểm chênh nhau bốn lần --- đó chính là toàn bộ giá trị của việc ưu
tiên kiểm thử.

\subsection{s33 · Dự phòng B2: vì sao chấm 30 lần --- {[}khi được
hỏi{]}}\label{s33-dux1ef1-phuxf2ng-b2-vuxec-sao-chux1ea5m-30-lux1ea7n-khi-ux111ux1b0ux1ee3c-hux1ecfi}

Chấm một lượt thì nhạy với thứ tự cố định của tập test, dễ may rủi. Nên
nhóm em chấm 30 lần, mỗi lần lấy ngẫu nhiên một mẫu con, rồi báo cáo
trung bình kèm độ lệch chuẩn. Với nhóm em, \textbf{độ lệch chuẩn nhiều
khi còn quan trọng hơn trung bình}, vì nó cho biết xếp hạng có ổn định
hay không.

\subsection{s34 · Dự phòng B3: bất biến lấy mẫu --- {[}khi được
hỏi{]}}\label{s34-dux1ef1-phuxf2ng-b3-bux1ea5t-biux1ebfn-lux1ea5y-mux1eabu-khi-ux111ux1b0ux1ee3c-hux1ecfi}

Nếu thầy hỏi về bất biến với tần số lấy mẫu: nhóm em lấy cùng một con
đường, resample ở nhiều mật độ điểm khác nhau rồi chấm lại. Độ chênh
khoảng 0,001 --- rất nhỏ nhưng \textbf{chưa bằng 0 như phép xoay}. Lý do
là cơ chế bias theo hiệu chiều dài cung mới chỉ gần bất biến tại tham số
hóa; đây đúng là lý do nhóm em muốn chuyển sang RoPE để đạt chính xác
tuyệt đối.

\subsection{s35 · Dự phòng B4: chi phí ở quy mô lớn --- {[}khi được
hỏi{]}}\label{s35-dux1ef1-phuxf2ng-b4-chi-phuxed-ux1edf-quy-muxf4-lux1edbn-khi-ux111ux1b0ux1ee3c-hux1ecfi}

Về chi phí: mô hình chỉ cần một lượt suy luận cho mỗi con đường, nên
chấm cả tập chỉ mất vài giây trên GPU. \textbf{Chi phí thật nằm ở mô
phỏng vật lý các test được chọn, không phải ở bộ chấm điểm.} Xếp hạng
tốt dồn lỗi lên đầu, nên chỉ cần chạy một phần nhỏ đầu hàng đợi đã lộ
phần lớn lỗi. Huấn luyện chỉ một lần 24 phút nên tái lập rất rẻ.

\subsection{s36 · Dự phòng B5: chứng minh ngắn tính bất biến --- {[}khi
được
hỏi{]}}\label{s36-dux1ef1-phuxf2ng-b5-chux1ee9ng-minh-ngux1eafn-tuxednh-bux1ea5t-biux1ebfn-khi-ux111ux1b0ux1ee3c-hux1ecfi}

Nếu thầy muốn chứng minh: dưới phép xoay-tịnh tiến, chiều dài mỗi đoạn
không đổi vì phép xoay bảo toàn khoảng cách và phép tịnh tiến thì triệt
tiêu; góc giữa hai đoạn kề cũng không đổi vì phép xoay bảo toàn góc; độ
cong và các đạo hàm theo chiều dài đều là hàm của các đại lượng bất biến
này nên cũng bất biến. Bỏ sin, cos, góc tuyệt đối là điều kiện cần. Hệ
quả: bảy kênh đầu vào giống hệt nhau sau biến đổi, nên mọi tầng phía sau
chỉ nhận đầu vào bất biến, và đầu ra bất biến.

\section{Chuẩn bị cho câu hỏi khó của
thầy}\label{chuux1ea9n-bux1ecb-cho-cuxe2u-hux1ecfi-khuxf3-cux1ee7a-thux1ea7y}

\subsection{Nếu thầy hỏi: ``SE2RoadNet APFD không cao hơn baseline, sao
gọi là đóng
góp?''}\label{nux1ebfu-thux1ea7y-hux1ecfi-se2roadnet-apfd-khuxf4ng-cao-hux1a1n-baseline-sao-gux1ecdi-luxe0-ux111uxf3ng-guxf3p}

Thưa thầy, nhóm em \textbf{không khẳng định thắng về APFD} --- APFD
ngang trong sai số. Đóng góp của nhóm em là \textbf{cải thiện Pareto}:
giữ APFD đỉnh, nâng AUC rõ rệt, và thêm một \textbf{bảo đảm bất biến
bằng 0} mà chưa phương pháp nào có. Trong bài toán an toàn, một bảo đảm
chứng minh được đáng giá hơn vài phần nghìn điểm số.

\subsection{Nếu thầy hỏi: ``Vì sao dám gọi là bất biến chính xác --- nó
là mạng nơ-ron
mà?''}\label{nux1ebfu-thux1ea7y-hux1ecfi-vuxec-sao-duxe1m-gux1ecdi-luxe0-bux1ea5t-biux1ebfn-chuxednh-xuxe1c-nuxf3-luxe0-mux1ea1ng-nux1a1-ron-muxe0}

Vì tính bất biến đến từ \textbf{khâu đặc trưng, không phải từ tham số
học được}. Bảy kênh không chứa tọa độ hay hướng tuyệt đối, nên xoay con
đường rồi trích lại cho đúng cùng một vector; mạng ở chế độ suy luận là
tất định, nên đầu ra giống hệt. Ở mức số thực có sai số cỡ một phần mười
triệu do ma trận xoay chứa sin/cos, nhưng ở độ chính xác báo cáo thì
APFD giống hệt đến từng bit.

\subsection{Nếu thầy hỏi: ``Vì sao tập Competition được coi là lệch phân
phối?''}\label{nux1ebfu-thux1ea7y-hux1ecfi-vuxec-sao-tux1eadp-competition-ux111ux1b0ux1ee3c-coi-luxe0-lux1ec7ch-phuxe2n-phux1ed1i}

Vì ba lý do: bộ sinh đường khác nhau, tỉ lệ FAIL khác nhau, và độ dài
đường khác hẳn --- Competition chỉ 129 tới 229 mét trong khi huấn luyện
dài tới hơn 450 mét. Bằng chứng là điểm trên tập kiểm tra cùng phân phối
và tập Competition chênh nhau đáng kể.

\subsection{Ghi chú nội bộ cho người trình bày (không
đọc)}\label{ghi-chuxfa-nux1ed9i-bux1ed9-cho-ngux1b0ux1eddi-truxecnh-buxe0y-khuxf4ng-ux111ux1ecdc}

Slide bước 4 hiện để Focal bằng 1,0 theo bản quét tham số mới nhất,
nhưng các slide headline và bảng xếp hạng vẫn đang hiển 0,8047 / 0,805 /
0,934 lấy từ lần chạy tham số 1,5. Hai con số này lệch nhau không đáng
kể và đều nằm trong sai số, nhưng nếu thầy soi kỹ, cứ trả lời thẳng: mọi
giá trị Focal đều cho APFD tương đương trong phạm vi độ lệch chuẩn, nên
kết luận không đổi.

\end{document}
